\begin{statementbox}
\textbf{统一约定。} $l$ 是切线，$n$ 是法线，角均由从 $P$ 出发的射线定义。
椭圆、双曲线中取非实轴顶点，使焦点三角形非退化；
抛物线中 $PH$ 平行对称轴且\textbf{指向开口方向}，取非顶点讨论非零角。
\begin{itemize}
\item 椭圆：\textbf{法线}包含 $\angle F_1PF_2$ 的内角平分射线，切线为外角平分线。
\item 双曲线：\textbf{切线}包含 $\angle F_1PF_2$ 的内角平分射线，法线为外角平分线。
\item 抛物线：\textbf{法线}包含 $\angle FPH$ 的内角平分射线，切线为外角平分线。
\end{itemize}
若将抛物线轴向射线改为指向开口的反方向，则内、外角平分线角色交换。
三类非退化曲线在顶点处也有唯一切线；这里只是排除退化的角或三角形。
\end{statementbox}
